

\section{The classifying space $\Gamma\backslash\left(\A_{\Q,S}\times\underline{E \Gamma}\right)$ and its codimension $1$ foliation} \label{sect4}

  Let $S=\{ p, q,\infty \}$ as above. The $C^*$-algebra $\cA=C_0(\A_{\Q,S})\rtimes \Gamma$ of the noncommutative space $X_{\Q,S}^{ab}:=\Gamma\backslash\A_{\Q,S}$ is not of type I (see \cite{CMbook}, Lemma 2.28 Chapter II).  The  understanding of $\cA$ at the level of $K$-theory is obtained by considering the assembly map \cite{BC} computing the $K$-theory of the  $C^*$-algebra $\cA$ in terms of  the universal  proper action  $\underline{E \Gamma}$ of $\Gamma$. Since  $\Gamma$ is commutative the various definitions of the cross product $C^*$-algebra all coincide with the reduced cross product, thus we drop the lower index $r$ in the notation $C^*_r$. Moreover the assembly map is an isomorphism. One obtains in this way  an ordinary space $Y=Y_{\Q,S}=\Gamma\backslash\left(\A_{\Q,S}\times\underline{E \Gamma}\right)$ which resolves the noncommutative nature of the quotient $\Gamma\backslash\A_{\Q,S}$. The group $\Gamma$ is $\Z^2\times\{ \pm 1\}$ and the universal  space for proper actions $\underline{E \Gamma}$ is $\R^2$ where $\Z^2$ acts by translations and $\{ \pm 1\}$ acts trivially. Thus one obtains  the following three dimensional space, quotient by the diagonal action of $\Gamma$,
 \begin{equation}\label{3space}
 Y=\Gamma\backslash\left(\A_{\Q,S}\times \underline{E \Gamma}\right) =\Gamma\backslash\left(\Q_p\times \Q_q\times \R \times \R^2\right).
 \end{equation}
 \vspace{.1in}
 
 We stress that the dimension three in the next proposition   is due to taking a pair of primes and has no relation to the homological dimension of $\Spec\,\Z$.
  \begin{prop}  	
 \label{bciso} $(i)$~The space $Y$ is locally compact of dimension $3$.\newline
  $(ii)$~The action of $\R^2$ by translations  on the   space $Y$ defines a codimension $1$  lamination of $Y$ whose space of leaves is $X_{\Q,S}^{ab}$.
    \end{prop}
    \proof  $(i)$~The properness of the action of $\Gamma$ on $\underline{E \Gamma}$ implies  that its action on $\left(\A_{\Q,S}\times \underline{E \Gamma}\right) $ is also proper. Thus $Y$ is Hausdorff  and since $\left(\A_{\Q,S}\times \underline{E \Gamma}\right) $ is locally compact, $Y$ is  locally compact. The dimension of $\Q_p\times \Q_q\times \R \times \R^2$ is equal to $3$ and is unaltered when passing to the quotient.\newline
    $(ii)$~By construction the orbits of the action by translation of $\R^2$ on $\left(\Q_p\times \Q_q\times \R \times \R^2\right)$ define a lamination which is invariant under the action of $\Gamma$ and hence descends to the proper quotient $Y$. The space of leaves of the first lamination is $\left(\Q_p\times \Q_q\times \R \right)$ and its quotient by $\Gamma$ is $X_{\Q,S}^{ab}$.\endproof 
     The simplest meaningful computation of the $K$-theory of the involved $C^*$-algebras  is for the cross product $A$ associated to the union in $X_{\Q,S}$ of the generic orbit with the three periodic orbits $C_p,C_q,C_\infty$. One obtains that $K_0(A)\simeq \Z^3$ reflects the presence of the three periodic orbits, while $K_1(A)\simeq \Z^2$ reflects the one-dimensionality of the periodic orbits $C_p,C_q$.
\begin{rem}    
$(i)$~The theory of the \'etale fundamental group for schemes can be viewed as the extension of Galois theory to schemes (\cite{Lenstra}). The classical abelian class field theory provides an isomorphism of the Galois group of the maximal abelian extension of a global field with an adelic group. It is natural to wonder if this %class field theory
 isomorphism can be extended to compare the behavior of the \'etale fundamental group of schemes with an adelic geometric construction.
Theorems \ref{mainglobalintro} and \ref{mainintro} %show that the class field theory isomorphism 
answer positively to this question for  schemes related to $\Spec\,\Z$ in terms of the geometry of the adele class space viewed as a covering $\pi:X_\Q^{ab}\to X_\Q$ of the scaling site. This extension of the class field theory isomorphism to schemes ought to play a role in the understanding of  the subtle relation between the scaling site and $\Spec\,\Z$.	\newline

$(ii)$~The scaling site which plays a central role in the present paper admits a natural complexification $\mathscr C_\Q$ \cite{CCsurvey}, whose role is to pass from characteristic one to characteristic zero. It replaces the real half-line by the pro-\'etale cover of the punctured unit disk in the complex plane. Then $\mathscr C_\Q$  is interpreted as the moduli space of elliptic curves with a triangular structure. The extrapolation of the findings of the present paper to the context of 
$\mathscr C_\Q$
  has the potential to formulate the results in   crystalline terms in place of the \'etale one. 
\end{rem}

